% ============================================================
\section*{Appendix}
\addcontentsline{toc}{section}{Appendix}

\appendix

\section{Quaternion Algebra}
    \label{app:quaternion}

    본 부록에서는 쿼터니언의 기본 대수 연산과 성질을 정리한다. 4장 및 5장의 수식 유도에 필요한 기초 내용을 포함한다.\\

    \subsection*{A.1 기본 정의}
        쿼터니언 $\mathbf{q} \in \mathbb{H}$는 실수 스칼라 성분 $q_0$와 허수 벡터 성분 $\mathbf{q}_v = [q_1, q_2, q_3]^\top$으로 구성된다:
 
        \begin{equation}
            \mathbf{q} = q_0 + q_1\mathbf{i} + q_2\mathbf{j} + q_3\mathbf{k} =
            \begin{bmatrix} q_0 \\ q_1 \\ q_2 \\ q_3 \end{bmatrix}
            \tag{A.1}
        \end{equation}
        
        허수 단위 사이의 곱셈 규칙:
        
        \begin{equation}
            \mathbf{i}^2 = \mathbf{j}^2 = \mathbf{k}^2 = \mathbf{ijk} = -1
            \tag{A.2}
        \end{equation}
        
        \begin{equation}
            \mathbf{ij} = \mathbf{k},\quad \mathbf{jk} = \mathbf{i},\quad \mathbf{ki} = \mathbf{j},\quad
            \mathbf{ji} = -\mathbf{k},\quad \mathbf{kj} = -\mathbf{i},\quad \mathbf{ik} = -\mathbf{j}
            \tag{A.3}
        \end{equation}
 
    \subsection*{A.2 기본 연산}
        \paragraph{노름 (Norm)}

            \begin{equation}
                \|\mathbf{q}\| = \sqrt{q_0^2 + q_1^2 + q_2^2 + q_3^2}
                \tag{A.4}
            \end{equation}
            
            단위 쿼터니언: $\|\mathbf{q}\| = 1$
 
        \paragraph{켤레 (Conjugate)}

            \begin{equation}
                \mathbf{q}^* = \begin{bmatrix} q_0 \\ -\mathbf{q}_v \end{bmatrix} =
                \begin{bmatrix} q_0 \\ -q_1 \\ -q_2 \\ -q_3 \end{bmatrix}
                \tag{A.5}
            \end{equation}
            
            단위 쿼터니언의 역원: $\mathbf{q}^{-1} = \mathbf{q}^*$

        \paragraph{쿼터니언 곱 (Hamilton Product)}
            두 쿼터니언 $\mathbf{p} = [p_0, \mathbf{p}_v]^\top$, $\mathbf{q} = [q_0, \mathbf{q}_v]^\top$의 곱:

            \begin{equation}
                \mathbf{p} \otimes \mathbf{q} =
                \begin{bmatrix}
                    p_0 q_0 - \mathbf{p}_v^\top \mathbf{q}_v \\
                    p_0 \mathbf{q}_v + q_0 \mathbf{p}_v + \mathbf{p}_v \times \mathbf{q}_v
                \end{bmatrix}
                \tag{A.6}
            \end{equation}
            
            쿼터니언 곱은 일반적으로 교환 법칙이 성립하지 않는다:
            $\mathbf{p} \otimes \mathbf{q} \neq \mathbf{q} \otimes \mathbf{p}$

        \paragraph{행렬 표현}
            좌측 곱 행렬 $\boldsymbol{\Xi}(\mathbf{p}) \in \mathbb{R}^{4\times 4}$:

            \begin{equation}
                \mathbf{p} \otimes \mathbf{q} = \boldsymbol{\Xi}(\mathbf{p})\mathbf{q} =
                \begin{bmatrix}
                    p_0 & -p_1 & -p_2 & -p_3 \\
                    p_1 &  p_0 & -p_3 &  p_2 \\
                    p_2 &  p_3 &  p_0 & -p_1 \\
                    p_3 & -p_2 &  p_1 &  p_0
                \end{bmatrix}
                \begin{bmatrix} q_0 \\ q_1 \\ q_2 \\ q_3 \end{bmatrix}
                \tag{A.7}
            \end{equation}
            
            우측 곱 행렬 $\boldsymbol{\Psi}(\mathbf{q}) \in \mathbb{R}^{4\times 4}$:
            
            \begin{equation}
                \mathbf{p} \otimes \mathbf{q} = \boldsymbol{\Psi}(\mathbf{q})\mathbf{p} =
                \begin{bmatrix}
                    q_0 & -q_1 & -q_2 & -q_3 \\
                    q_1 &  q_0 &  q_3 & -q_2 \\
                    q_2 & -q_3 &  q_0 &  q_1 \\
                    q_3 &  q_2 & -q_1 &  q_0
                \end{bmatrix}
                \begin{bmatrix} p_0 \\ p_1 \\ p_2 \\ p_3 \end{bmatrix}
                \tag{A.8}
            \end{equation}
 
    \subsection*{A.3 회전 표현}
 
회전 축 $\hat{\mathbf{u}} \in \mathbb{R}^3$ ($\|\hat{\mathbf{u}}\| = 1$)에 대한
각도 $\alpha$ 회전의 단위 쿼터니언:
 
\begin{equation}
    \mathbf{q} = \begin{bmatrix} \cos(\alpha/2) \\ \hat{\mathbf{u}}\sin(\alpha/2) \end{bmatrix}
    \tag{A.9}
\end{equation}
 
벡터 $\mathbf{v} \in \mathbb{R}^3$의 쿼터니언 회전:
 
\begin{equation}
    \begin{bmatrix} 0 \\ \mathbf{v}' \end{bmatrix} =
    \mathbf{q} \otimes \begin{bmatrix} 0 \\ \mathbf{v} \end{bmatrix} \otimes \mathbf{q}^*
    \tag{A.10}
\end{equation}
 
쿼터니언-DCM 변환:
 
\begin{equation}
    \mathbf{R}(\mathbf{q}) =
    \begin{bmatrix}
        1-2(q_2^2+q_3^2) & 2(q_1q_2+q_0q_3) & 2(q_1q_3-q_0q_2) \\
        2(q_1q_2-q_0q_3) & 1-2(q_1^2+q_3^2) & 2(q_2q_3+q_0q_1) \\
        2(q_1q_3+q_0q_2) & 2(q_2q_3-q_0q_1) & 1-2(q_1^2+q_2^2)
    \end{bmatrix}
    \tag{A.11}
\end{equation}
 
% ------------------------------------------------------------
\subsection*{A.4 쿼터니언 보간 (SLERP)}
 
두 단위 쿼터니언 $\mathbf{q}_0$, $\mathbf{q}_1$ 사이의
구면 선형 보간(Spherical Linear Interpolation, SLERP):
 
\begin{equation}
    \text{Slerp}(\mathbf{q}_0, \mathbf{q}_1; t) =
    \mathbf{q}_0 \left(\mathbf{q}_0^* \otimes \mathbf{q}_1\right)^t, \quad t \in [0, 1]
    \tag{A.12}
\end{equation}
 
닫힌 형태:
 
\begin{equation}
    \text{Slerp}(\mathbf{q}_0, \mathbf{q}_1; t) =
    \frac{\sin((1-t)\Omega)}{\sin\Omega}\mathbf{q}_0 +
    \frac{\sin(t\Omega)}{\sin\Omega}\mathbf{q}_1, \quad
    \cos\Omega = \mathbf{q}_0^\top \mathbf{q}_1
    \tag{A.13}
\end{equation}
 
% ============================================================
\section{EKF Jacobian Derivation}
\label{app:jacobian}
% ============================================================
 
본 부록에서는 5장 EKF 설계에 사용된 Jacobian 행렬의
전체 유도 과정을 기술한다.
 
% ------------------------------------------------------------
\subsection*{B.1 상태 방정식 Jacobian $\mathbf{F}_c$}
 
상태 벡터 $\mathbf{x} = [\mathbf{q}^\top, \mathbf{b}_g^\top]^\top \in \mathbb{R}^7$에 대한
연속시간 상태 방정식:
 
\begin{equation}
    \dot{\mathbf{x}} = \mathbf{f}(\mathbf{x}, \tilde{\boldsymbol{\omega}}) =
    \begin{bmatrix}
        \frac{1}{2}\boldsymbol{\Omega}(\tilde{\boldsymbol{\omega}} - \mathbf{b}_g)\mathbf{q} \\
        \mathbf{0}_3
    \end{bmatrix}
    \tag{B.1}
\end{equation}
 
\paragraph{$\partial \dot{\mathbf{q}} / \partial \mathbf{q}$ 계산}
 
$\dot{\mathbf{q}} = \frac{1}{2}\boldsymbol{\Omega}(\hat{\boldsymbol{\omega}})\mathbf{q}$이므로:
 
\begin{equation}
    \frac{\partial \dot{\mathbf{q}}}{\partial \mathbf{q}} =
    \frac{1}{2}\boldsymbol{\Omega}(\hat{\boldsymbol{\omega}}) =
    \frac{1}{2}\begin{bmatrix}
         0       & -\hat{\omega}_x & -\hat{\omega}_y & -\hat{\omega}_z \\
         \hat{\omega}_x &  0        &  \hat{\omega}_z & -\hat{\omega}_y \\
         \hat{\omega}_y & -\hat{\omega}_z &  0        &  \hat{\omega}_x \\
         \hat{\omega}_z &  \hat{\omega}_y & -\hat{\omega}_x &  0
    \end{bmatrix}
    \in \mathbb{R}^{4\times 4}
    \tag{B.2}
\end{equation}
 
\noindent 여기서 $\hat{\boldsymbol{\omega}} = \tilde{\boldsymbol{\omega}} - \hat{\mathbf{b}}_g$이다.
 
\paragraph{$\partial \dot{\mathbf{q}} / \partial \mathbf{b}_g$ 계산}
 
$\dot{\mathbf{q}} = \frac{1}{2}\boldsymbol{\Xi}(\mathbf{q})\begin{bmatrix}0\\\hat{\boldsymbol{\omega}}\end{bmatrix}$이고,
$\hat{\boldsymbol{\omega}} = \tilde{\boldsymbol{\omega}} - \mathbf{b}_g$이므로:
 
\begin{equation}
    \frac{\partial \dot{\mathbf{q}}}{\partial \mathbf{b}_g} =
    -\frac{1}{2}\boldsymbol{\Xi}_{1:3}(\mathbf{q}) \in \mathbb{R}^{4\times 3}
    \tag{B.3}
\end{equation}
 
\noindent 여기서 $\boldsymbol{\Xi}_{1:3}(\mathbf{q})$는 $\boldsymbol{\Xi}(\mathbf{q})$의
우측 3열로 구성된 부분 행렬:
 
\begin{equation}
    \boldsymbol{\Xi}_{1:3}(\mathbf{q}) =
    \begin{bmatrix}
        -q_1 & -q_2 & -q_3 \\
         q_0 & -q_3 &  q_2 \\
         q_3 &  q_0 & -q_1 \\
        -q_2 &  q_1 &  q_0
    \end{bmatrix}
    \tag{B.4}
\end{equation}
 
\paragraph{전체 $\mathbf{F}_c$}
 
\begin{equation}
    \mathbf{F}_c = \frac{\partial \mathbf{f}}{\partial \mathbf{x}} =
    \begin{bmatrix}
        \frac{1}{2}\boldsymbol{\Omega}(\hat{\boldsymbol{\omega}}) &
        -\frac{1}{2}\boldsymbol{\Xi}_{1:3}(\hat{\mathbf{q}}) \\[6pt]
        \mathbf{0}_{3\times 4} & \mathbf{0}_{3\times 3}
    \end{bmatrix}
    \in \mathbb{R}^{7\times 7}
    \tag{B.5}
\end{equation}
 
이산시간 상태 천이 행렬:
 
\begin{equation}
    \mathbf{F}_k = \mathbf{I}_7 + \mathbf{F}_c\,\Delta t
    \tag{B.6}
\end{equation}
 
% ------------------------------------------------------------
\subsection*{B.2 가속도계 관측 Jacobian $\mathbf{H}_a$}
 
관측 함수 $\mathbf{h}_a(\mathbf{q}) = \mathbf{R}(\mathbf{q})\mathbf{g}^n$,
$\mathbf{g}^n = [0, 0, g]^\top$을 전개하면:
 
\begin{equation}
    \mathbf{h}_a(\mathbf{q}) = g\begin{bmatrix}
        2(q_1 q_3 - q_0 q_2) \\
        2(q_2 q_3 + q_0 q_1) \\
        q_0^2 - q_1^2 - q_2^2 + q_3^2
    \end{bmatrix}
    \tag{B.7}
\end{equation}
 
각 성분을 $q_0, q_1, q_2, q_3$에 대해 편미분하면:
 
\begin{align}
    \frac{\partial h_{a,1}}{\partial \mathbf{q}} &= 2g\begin{bmatrix} -q_2 & q_3 & -q_0 & q_1 \end{bmatrix} \tag{B.8} \\
    \frac{\partial h_{a,2}}{\partial \mathbf{q}} &= 2g\begin{bmatrix}  q_1 & q_0 &  q_3 & q_2 \end{bmatrix} \tag{B.9} \\
    \frac{\partial h_{a,3}}{\partial \mathbf{q}} &= 2g\begin{bmatrix}  q_0 & -q_1 & -q_2 & q_3 \end{bmatrix} \tag{B.10}
\end{align}
 
따라서:
 
\begin{equation}
    \frac{\partial \mathbf{h}_a}{\partial \mathbf{q}} = 2g
    \begin{bmatrix}
        -q_2 &  q_3 & -q_0 &  q_1 \\
         q_1 &  q_0 &  q_3 &  q_2 \\
         q_0 & -q_1 & -q_2 &  q_3
    \end{bmatrix}
    \in \mathbb{R}^{3\times 4}
    \tag{B.11}
\end{equation}
 
바이어스에 대한 편미분은 $\mathbf{0}_{3\times 3}$이므로
전체 관측 Jacobian:
 
\begin{equation}
    \mathbf{H}_{a} = \left[\frac{\partial \mathbf{h}_a}{\partial \mathbf{q}} \;\Bigg|\; \mathbf{0}_{3\times 3}\right]
    \in \mathbb{R}^{3\times 7}
    \tag{B.12}
\end{equation}
 
% ------------------------------------------------------------
\subsection*{B.3 지자기계 관측 Jacobian $\mathbf{H}_m$}
 
관측 함수 $\mathbf{h}_m(\mathbf{q}) = \mathbf{R}(\mathbf{q})\mathbf{m}^n$,
$\mathbf{m}^n = [m_N, 0, m_D]^\top$을 전개하면:
 
\begin{equation}
    \mathbf{h}_m(\mathbf{q}) =
    \begin{bmatrix}
        m_N(1 - 2q_2^2 - 2q_3^2) + 2m_D(q_1 q_3 - q_0 q_2) \\
        2m_N(q_1 q_2 + q_0 q_3) + 2m_D(q_2 q_3 - q_0 q_1) \\
        2m_N(q_1 q_3 - q_0 q_2) + m_D(1 - 2q_1^2 - 2q_2^2)
    \end{bmatrix}
    \tag{B.13}
\end{equation}
 
각 성분을 $q_0, q_1, q_2, q_3$에 대해 편미분하면:
 
\begin{equation}
    \frac{\partial \mathbf{h}_m}{\partial \mathbf{q}} = 2
    \begin{bmatrix}
        -m_D q_2 & m_D q_3 & -2m_N q_2 - m_D q_0 & -2m_N q_3 + m_D q_1 \\
         m_N q_3 - m_D q_1 & m_N q_2 + m_D q_0 & m_N q_1 - m_D q_3 & m_N q_0 + m_D q_2 \\
        -m_N q_2 & m_N q_3 - 2m_D q_1 & -m_N q_0 - 2m_D q_2 & m_N q_1
    \end{bmatrix}
    \tag{B.14}
\end{equation}
 
전체 관측 Jacobian:
 
\begin{equation}
    \mathbf{H}_{m} = \left[\frac{\partial \mathbf{h}_m}{\partial \mathbf{q}} \;\Bigg|\; \mathbf{0}_{3\times 3}\right]
    \in \mathbb{R}^{3\times 7}
    \tag{B.15}
\end{equation}
 
% ============================================================
\section{Code Structure}
\label{app:code}
% ============================================================
 
본 부록에서는 Python으로 구현한 시뮬레이션 및 AHRS 알고리즘의
코드 구조를 설명한다.
 
% ------------------------------------------------------------
\subsection*{C.1 전체 디렉토리 구조}
 
\begin{verbatim}
prj_imu_fusion/
├── sim/
│   ├── imu_simulator.py      # IMU 데이터 합성
│   ├── trajectory.py         # 참 자세 궤적 생성
│   └── noise_model.py        # 잡음 모델 (ARW, RRW)
├── filters/
│   ├── complementary.py      # Complementary Filter
│   ├── mahony.py             # Mahony Filter
│   ├── madgwick.py           # Madgwick Filter
│   └── ekf.py                # Extended Kalman Filter
├── calibration/
│   ├── accel_calib.py        # 가속도계 캘리브레이션
│   ├── gyro_calib.py         # 자이로 Allan Variance
│   └── mag_calib.py          # Hard/Soft Iron 보정
├── analysis/
│   ├── allan_variance.py     # AVAR 계산 및 시각화
│   └── psd.py                # PSD 분석
├── hardware/
│   └── mpu9250_reader.py     # Arduino 시리얼 통신
└── main.py                   # 메인 실행 스크립트
\end{verbatim}
 
% ------------------------------------------------------------
\subsection*{C.2 핵심 모듈 설명}
 
\paragraph{imu\_simulator.py}
 
\begin{verbatim}
class IMUSimulator:
    def __init__(self, sigma_w, sigma_rw, sigma_a, sigma_m, fs):
        # 잡음 파라미터 초기화
    
    def generate(self, q_true, omega_true):
        # 자이로, 가속도계, 지자기계 측정값 합성
        # 반환: gyro_meas, accel_meas, mag_meas
\end{verbatim}
 
\paragraph{ekf.py}
 
\begin{verbatim}
class EKF:
    def __init__(self, Q, Ra, Rm, P0):
        # 공분산 초기화
        self.x = np.array([1,0,0,0, 0,0,0])  # 상태 벡터
        self.P = P0
    
    def predict(self, gyro, dt):
        # 상태 전파 및 공분산 예측
        # F = I + Fc * dt
    
    def update_accel(self, accel):
        # 가속도계 Update (동적 가속도 조건 포함)
        # H_a, K_a, 상태/공분산 갱신
    
    def update_mag(self, mag):
        # 지자기계 Update
        # H_m, K_m, 상태/공분산 갱신
    
    def normalize(self):
        # 쿼터니언 정규화
    
    def get_euler(self):
        # 쿼터니언 → 오일러각 변환
\end{verbatim}
 
\paragraph{allan\_variance.py}
 
\begin{verbatim}
def compute_avar(data, fs):
    # 입력: 정지 상태 자이로 시계열, 샘플링 주파수
    # 출력: tau 배열, AVAR 배열
    # 방법: overlapping AVAR
 
def extract_noise_params(tau, avar):
    # log-log 기울기 분석으로 ARW, BI, RRW 계수 추출
    # 반환: N (ARW), B (BI), K (RRW)
\end{verbatim}
 
% ------------------------------------------------------------
\subsection*{C.3 실행 방법}
 
\paragraph{시뮬레이션 실행}
 
\begin{verbatim}
# 의존성 설치
pip install numpy scipy matplotlib
 
# 시뮬레이션 실행
python main.py --mode sim --filter ekf --duration 60
 
# 필터 비교
python main.py --mode sim --filter all --plot
\end{verbatim}
 
\paragraph{실제 데이터 분석}
 
\begin{verbatim}
# Arduino 시리얼 데이터 수집 (baudrate: 115200)
python hardware/mpu9250_reader.py --port COM3 --duration 7200
 
# Allan Variance 분석
python analysis/allan_variance.py --input data/gyro_static.csv
 
# 캘리브레이션 실행
python calibration/accel_calib.py --input data/accel_6pos.csv
python calibration/mag_calib.py   --input data/mag_rotation.csv
\end{verbatim}
 
% ============================================================
\section{MPU-9250 Hardware Overview}
\label{app:hardware}
% ============================================================
 
본 부록에서는 MPU-9250 IMU의 하드웨어 사양과
통신 인터페이스를 정리한다.
 
% ------------------------------------------------------------
\subsection*{D.1 센서 사양}
 
\begin{table}[H]
    \centering
    \caption{MPU-9250 센서 사양}
    \label{tab:mpu9250_spec}
    \begin{tabular}{llll}
        \toprule
        항목 & 가속도계 & 자이로스코프 & 지자기계 (AK8963) \\
        \midrule
        측정 범위 & $\pm$2/4/8/16 g & $\pm$250/500/1000/2000 $^\circ$/s & $\pm$4800 $\mu$T \\
        ADC 해상도 & 16 bit & 16 bit & 14 / 16 bit \\
        감도 (기본) & 16384 LSB/g & 131 LSB/($^\circ$/s) & 0.6 $\mu$T/LSB \\
        출력 데이터 레이트 & 최대 4 kHz & 최대 8 kHz & 최대 100 Hz \\
        동작 전압 & \multicolumn{3}{c}{1.71 -- 3.45 V} \\
        소비 전류 & 450 $\mu$A & 3.2 mA & 280 $\mu$A \\
        온도 범위 & \multicolumn{3}{c}{$-40^\circ$C $\sim$ $+85^\circ$C} \\
        패키지 & \multicolumn{3}{c}{QFN 3$\times$3$\times$1 mm} \\
        \bottomrule
    \end{tabular}
\end{table}
 
% ------------------------------------------------------------
\subsection*{D.2 I$^2$C 인터페이스}
 
MPU-9250은 호스트와 I$^2$C 프로토콜로 통신하며,
AD0 핀 레벨에 따라 슬레이브 주소가 결정된다:
 
\begin{equation}
    \text{AD0} = \text{LOW} \Rightarrow 0\text{x}68, \qquad
    \text{AD0} = \text{HIGH} \Rightarrow 0\text{x}69
    \tag{D.1}
\end{equation}
 
주요 레지스터:
 
\begin{table}[H]
    \centering
    \caption{MPU-9250 주요 레지스터}
    \label{tab:mpu9250_reg}
    \begin{tabular}{lll}
        \toprule
        레지스터 주소 & 이름 & 설명 \\
        \midrule
        0x3B -- 0x40 & ACCEL\_XOUT -- ACCEL\_ZOUT & 가속도계 원시 데이터 (16비트 $\times$ 3축) \\
        0x41 -- 0x42 & TEMP\_OUT & 온도 데이터 \\
        0x43 -- 0x48 & GYRO\_XOUT -- GYRO\_ZOUT & 자이로 원시 데이터 (16비트 $\times$ 3축) \\
        0x6B & PWR\_MGMT\_1 & 전원 관리 및 클럭 설정 \\
        0x1B & GYRO\_CONFIG & 자이로 측정 범위 설정 \\
        0x1C & ACCEL\_CONFIG & 가속도계 측정 범위 설정 \\
        0x37 & INT\_PIN\_CFG & Pass-Through 모드 활성화 \\
        \bottomrule
    \end{tabular}
\end{table}
 
원시 데이터 변환:
 
\begin{equation}
    \text{raw} = (\text{HIGH\_BYTE} \ll 8) \mid \text{LOW\_BYTE}
    \tag{D.2}
\end{equation}
 
\begin{equation}
    \text{physical value} = \frac{\text{raw}}{\text{sensitivity}}
    \tag{D.3}
\end{equation}
 
% ------------------------------------------------------------
\subsection*{D.3 AK8963 지자기계 접근}
 
AK8963은 MPU-9250의 보조 I$^2$C 버스에 연결되어 있으며,
Pass-Through 모드를 활성화하면 호스트에서 직접 접근할 수 있다.
 
\begin{enumerate}
    \item MPU-9250 레지스터 0x6A (USER\_CTRL)의 I2C\_MST\_EN 비트를 0으로 설정
    \item MPU-9250 레지스터 0x37 (INT\_PIN\_CFG)의 BYPASS\_EN 비트를 1로 설정
    \item AK8963 슬레이브 주소 0x0C로 직접 접근
    \item AK8963 레지스터 0x0A (CNTL1)에 0x16 기록: 16비트 연속 측정 모드 2
\end{enumerate}
 
AK8963 주요 레지스터:
 
\begin{table}[H]
    \centering
    \caption{AK8963 주요 레지스터}
    \label{tab:ak8963_reg}
    \begin{tabular}{lll}
        \toprule
        레지스터 주소 & 이름 & 설명 \\
        \midrule
        0x02 & ST1 & 데이터 준비 상태 (DRDY 비트) \\
        0x03 -- 0x08 & HXL -- HZH & 자기장 원시 데이터 (16비트 $\times$ 3축) \\
        0x09 & ST2 & 데이터 오버런 및 자기 센서 오버플로 \\
        0x0A & CNTL1 & 측정 모드 및 출력 비트 설정 \\
        0x10 -- 0x12 & ASAX/Y/Z & 감도 보정값 (ROM) \\
        \bottomrule
    \end{tabular}
\end{table}
 
감도 보정값(Sensitivity Adjustment Value) 적용:
 
\begin{equation}
    H_{\text{adj},i} = H_i \times \frac{(ASA_i - 128) \times 0.5}{128} + H_i, \quad i \in \{x, y, z\}
    \tag{D.4}
\end{equation}
 
% ------------------------------------------------------------
\subsection*{D.4 좌표계 정렬}
 
MPU-9250 가속도계/자이로스코프와 AK8963 지자기계는
서로 다른 좌표계를 사용하므로 데이터 퓨전 전 축 변환이 필요하다.
 
\begin{table}[H]
    \centering
    \caption{MPU-9250 좌표계 정렬}
    \label{tab:axis_align}
    \begin{tabular}{lccc}
        \toprule
        센서 & X축 & Y축 & Z축 \\
        \midrule
        가속도계 / 자이로스코프 & 전방 & 우측 & 하방 \\
        AK8963 지자기계 & 우측 & 전방 & 하방 \\
        \midrule
        정렬 후 (Body Frame) & 전방 & 우측 & 하방 \\
        \bottomrule
    \end{tabular}
\end{table}
 
AK8963 출력의 Body Frame 변환:
 
\begin{equation}
    \begin{bmatrix} m_x^b \\ m_y^b \\ m_z^b \end{bmatrix} =
    \begin{bmatrix} m_y^{\text{AK}} \\ m_x^{\text{AK}} \\ -m_z^{\text{AK}} \end{bmatrix}
    \tag{D.5}
\end{equation}